\section{Язык задачи}
\label{sec:setup}

В этой главе введены переменные и гамильтониан, на которых построено всё дальнейшее. Независимая переменная --- длина $s$ вдоль референсной орбиты, поперечные импульсы нормированы на импульс синхронной частицы. Отклонение энергии задаётся тремя способами --- через импульс, полную и кинетическую энергию; приведены формулы перевода хроматичности между ними, поскольку разные программы пользуются разными переменными.

\subsection{Обозначения}

Индекс $s$ отмечает синхронную (референсную) частицу.

\begin{center}
\renewcommand{\arraystretch}{1.2}
\begin{tabular}{@{}lp{0.55\textwidth}@{}}
\hline
\multicolumn{2}{@{}l}{\emph{Синхронная частица}}\\
$p_s$, $E_s=\gamma_s mc^2$, $K_s$ & импульс, полная и кинетическая энергия\\
$\beta_s$, $\gamma_s$ & скорость в единицах $c$ и лоренц-фактор\\
$C_s$ & длина орбиты\\
\hline
\multicolumn{2}{@{}l}{\emph{Оптика}}\\
$1/\rho$ & кривизна референсной орбиты (в диполе), $0$ вне диполей\\
$K_1$, $K_2$ & градиент квадруполя и секступоля, нормированные на $B\rho=p_s/q$\\
$K_x=K_1+1/\rho^2$, $K_y=-K_1$ & жёсткости линейных уравнений\\
$\alpha$, $\beta$, $\gamma$ & функции Твисса (по $x$ или $y$)\\
$D_1$, $D_1'$ & дисперсия на $\delta$ и её производная\\
$\nu_{x,y}$, $\mu=2\pi\nu$ & числа бетатронных колебаний и набег фазы за оборот\\
$J$, $\psi$ & инвариант Куранта--Снайдера и бетатронная фаза\\
$\langle\cdot\rangle_\psi$ & среднее по бетатронной фазе\\
\hline
\multicolumn{2}{@{}l}{\emph{Отклонение энергии}}\\
$\delta$ & $=\delta_p$, относительное отклонение импульса \eqref{eq:dp}\\
$\delta_E$, $\delta_K$ & то же для полной и кинетической энергии \eqref{eq:dE-dK}\\
$\xi=\mathrm{d}\nu/\mathrm{d}\delta$ & хроматичность\\
\hline
\multicolumn{2}{@{}l}{\emph{Резонатор}}\\
$h$, $\omega_\rf$, $V$ & кратность, частота, амплитуда напряжения\\
\hline
\end{tabular}
\end{center}

Буква $h$ закреплена за кратностью ВЧ; кривизна всюду пишется как $1/\rho$.

\subsection{Гамильтониан в сопровождающей системе}

Сопровождающая система координат связана с референсной орбитой: $x$ --- отклонение в плоскости поворота (наружу от центра кривизны), $y$ --- вертикальное. Канонические импульсы $p_x$, $p_y$ нормированы на $p_s$; вне векторного потенциала $p_x=P_x/p_s$, где $P_x$ --- механический импульс. Гамильтониан есть взятая с минусом $s$-компонента канонического импульса, делённая на $p_s$. Элемент дуги на расстоянии $x$ от орбиты длиннее $\mathrm{d}s$ в $1+x/\rho$ раз, поэтому кинетическая часть несёт этот множитель. В диполе с однородным полем $B_0$, $qB_0/p_s=1/\rho$, векторный потенциал в сопровождающих координатах сводится к полиному. Вместе с полями квадруполя и секступоля
\begin{equation}
\label{eq:H}
\mathcal{H}
=-\Bigl(1+\frac{x}{\rho}\Bigr)\sqrt{(1+\delta)^2-p_x^2-p_y^2}
+\frac{x}{\rho}+\frac{x^2}{2\rho^2}
+\frac{K_1}{2}\bigl(x^2-y^2\bigr)
+\frac{K_2}{6}\bigl(x^3-3xy^2\bigr).
\end{equation}
Поля в \eqref{eq:H} фиксированы: от $\delta$ зависит только кинетическая часть. Это главное свойство канонической формулировки --- вся хроматичность происходит из корня и множителя $1+x/\rho$.

В линейном порядке уравнения Гамильтона дают
\begin{equation}
\label{eq:hill}
x''+\frac{K_x}{1+\delta}\,x=\frac{\delta}{\rho(1+\delta)},
\qquad
y''+\frac{K_y}{1+\delta}\,y=0,
\end{equation}
т.\,е. при импульсе $p_s(1+\delta)$ жёсткости делятся на $1+\delta$.

Если ВЧ нет, $\delta$ --- сохраняющийся параметр, и движение четырёхмерно. С ВЧ добавляется пара (продольная координата, $\delta$); в COSY это $(l,\delta_K)$ с $l=-(t-t_0)v_0\gamma_s/(1+\gamma_s)$~\cite{ValetovPhD}. В статическом поле вместо $\delta$ можно взять пару $(t,-E)$, как в~\cite{EqNote}. До главы~\ref{sec:rf} ВЧ нет.

\subsection{Переменные энергии и перевод хроматичности}

Определим
\begin{equation}
\label{eq:dp}
\delta=\delta_p=\frac{p-p_s}{p_s},
\end{equation}
\begin{equation}
\label{eq:dE-dK}
\delta_E=\frac{E-E_s}{E_s},
\qquad
\delta_K=\frac{K-K_s}{K_s},
\qquad
p_t=\frac{E-E_s}{p_sc}.
\end{equation}
Последняя переменная --- \texttt{PT} в MAD-X~\cite{MADX}. Из $\mathrm{d}E=v\,\mathrm{d}p$ и $\mathrm{d}K=\mathrm{d}E$ в первом порядке
\begin{equation}
\label{eq:conv}
\delta_E=\beta_s^2\,\delta,
\qquad
p_t=\beta_s\,\delta,
\qquad
\delta_K=\Bigl(1+\frac{1}{\gamma_s}\Bigr)\delta,
\end{equation}
где использовано $E_s\beta_s^2/K_s=(\gamma_s^2-1)/(\gamma_s(\gamma_s-1))=(\gamma_s+1)/\gamma_s$. Хроматичность --- производная, поэтому
\begin{equation}
\label{eq:xi-conv}
\xi\equiv\frac{\mathrm{d}\nu}{\mathrm{d}\delta}
=\Bigl(1+\frac{1}{\gamma_s}\Bigr)\frac{\mathrm{d}\nu}{\mathrm{d}\delta_K}
=\beta_s\frac{\mathrm{d}\nu}{\mathrm{d}p_t}
=\beta_s^2\frac{\mathrm{d}\nu}{\mathrm{d}\delta_E}.
\end{equation}
Перевод точен для первой производной в точке $\delta=0$. Для нелинейной хроматичности \eqref{eq:conv} нужно брать с поправками следующих порядков (например, $\delta_K$ содержит член $\propto\delta^2$), и коэффициенты при $\delta^j$, $j\ge2$, переводятся не простым умножением. Ниже везде $\xi=\mathrm{d}\nu/\mathrm{d}\delta$ при $\delta=0$ и бесконечно малой амплитуде, натуральная, если не оговорено иное.
